\documentclass[11pt,a4paper]{article}
\usepackage[utf8]{inputenc}
\usepackage[T1]{fontenc}
\usepackage[margin=1in]{geometry}
\usepackage{amsmath,amssymb}
\usepackage{booktabs}
\usepackage{multirow}
\usepackage{array}
\usepackage{caption}
\usepackage{float}
\usepackage{hyperref}
\usepackage[table,xcdraw]{xcolor}
\usepackage{listings}
\usepackage{parskip}
\usepackage{adjustbox}

\hypersetup{colorlinks=true,linkcolor=blue,urlcolor=blue}

\definecolor{lightgray}{gray}{0.93}
\definecolor{bestgreen}{rgb}{0.85,1.0,0.85}
\definecolor{badred}{rgb}{1.0,0.82,0.82}
\definecolor{midyellow}{rgb}{1.0,0.97,0.80}
\lstset{basicstyle=\small\ttfamily,frame=single,backgroundcolor=\color{lightgray},breaklines=true}

\title{\textbf{LLM-Based Financial Statement Forecasting} \\[4pt]
       \large Part~2: Gemini 2.5~Flash, Ensemble Methods, \\
       Robustness Testing, and Thinking-Mode Analysis}
\author{Feng Qiu and Alice (Ying) Xia}
\date{February 2026}

\begin{document}
\maketitle
\tableofcontents
\newpage

%───────────────────────────────────────────────────────────────────────────────
\section{Introduction}
%───────────────────────────────────────────────────────────────────────────────

This report documents the complete Part~2 research pipeline evaluating
Google Gemini~2.5~Flash for financial statement forecasting in a banking
lending context.  We address five questions:

\begin{enumerate}
  \item Does Gemini match or outperform a policy-based model on 50~S\&P~500 companies?
  \item Can ensemble combinations reduce forecasting error beyond either model alone?
  \item How does temperature affect both accuracy and output consistency?
  \item Does Gemini's \emph{thinking} mode improve accuracy?
  \item Does thinking mode hurt or help forecast \emph{consistency} across 10 repeated runs?
\end{enumerate}

Questions 4 and~5 are deliberately kept separate: accuracy is about
\emph{proximity to actual values}, while robustness (consistency) is about
\emph{reproducibility across runs with the same input}.

%───────────────────────────────────────────────────────────────────────────────
\section{Data and Metrics}
%───────────────────────────────────────────────────────────────────────────────

\paragraph{Dataset.}
50~S\&P~500 companies, annual financial statements 2020--2023 as input,
2024 actual statements (from SEC~EDGAR) as ground truth.

\paragraph{Accuracy metric.}  Symmetric Mean Absolute Percentage Error (sMAPE):
\begin{equation}
  \text{sMAPE} = \frac{100}{n}\sum_{i=1}^{n}
    \frac{\lvert F_i-A_i\rvert}{\tfrac{1}{2}(\lvert F_i\rvert+\lvert A_i\rvert)}
  \label{eq:smape}
\end{equation}

\paragraph{Robustness metric.}  Coefficient of Variation (CV) across $n$ repeated runs:
\begin{equation}
  \text{CV} = \frac{\sigma}{\lvert\mu\rvert}\times100\%
  \label{eq:cv}
\end{equation}
CV measures \emph{how much the model's output varies} when given
identical inputs.  It has no relation to accuracy.
A perfectly consistent model has CV~$= 0\%$ regardless of whether its
forecast is close to or far from the actual value.

%───────────────────────────────────────────────────────────────────────────────
\section{LLM Forecasting: Gemini~2.5~Flash}
%───────────────────────────────────────────────────────────────────────────────

\subsection{Model Configuration}

\begin{table}[H]\centering
\caption{Gemini~2.5~Flash production baseline configuration}
\label{tab:config}
\begin{adjustbox}{max width=\textwidth}
\begin{tabular}{ll}\toprule
\textbf{Parameter} & \textbf{Value}\\\midrule
Model ID            & \texttt{gemini-2.5-flash}\\
Temperature         & 0.0 (deterministic)\\
Top-p               & 0.95\\
Max output tokens   & 8{,}192 (optimised from initial 4{,}096)\\
Response format     & JSON mode (\texttt{application/json})\\
Rate limit          & 15 requests / minute\\
Thinking budget     & Not set (adaptive default)\\
\bottomrule
\end{tabular}
\end{adjustbox}
\end{table}

\subsection{Prompt Design}

A compact prompt supplies four years of key financials and requests a
fixed-schema JSON forecast:

\begin{lstlisting}[caption={Forecast prompt (condensed)}]
Forecast 2024 financials for {TICKER} based on 2020-2023 data.
Historical Data:
  2020: Rev $X, NI $Y, Assets $Z, Equity $W
  ...
Task: Analyze trends and forecast 2024.
Output ONLY this JSON (NUMBERS ONLY, NO TEXT):
{ "forecast_2024": { "income_statement": {...},
                     "balance_sheet": {...},
                     "cash_flow": {...} } }
\end{lstlisting}

\subsection{Error Handling}

\begin{enumerate}
  \item \textbf{Retry} -- up to 3 attempts; temperature steps 0.0 $\to$ 0.3 on retry.
  \item \textbf{LLM JSON repair} -- secondary call extracts numbers from truncated JSON.
  \item \textbf{Timeout} -- 60-second hard limit via \texttt{SIGALRM}.
  \item \textbf{Graceful degradation} -- fall back to traditional forecast on total failure.
\end{enumerate}

\subsection{Accuracy: Gemini vs.\ Traditional}

\begin{table}[H]\centering
\caption{Gemini~2.5~Flash vs.\ Traditional Model -- 50 S\&P~500 companies}
\label{tab:model_comparison}
\begin{adjustbox}{max width=\textwidth}
\begin{tabular}{lcc}\toprule
\textbf{Metric} & \textbf{Gemini 2.5 Flash} & \textbf{Traditional (V\'elez-Pareja)}\\\midrule
Mean sMAPE   & 43.46\% & 36.53\%\\
Median sMAPE & 29.93\% & 30.80\%\\
Std Dev      & 45.82\% & 28.47\%\\\midrule
\multicolumn{3}{l}{\textit{Head-to-head (lower sMAPE wins)}}\\
Gemini wins      & \multicolumn{2}{c}{26/50 \quad(52.0\%)}\\
Traditional wins & \multicolumn{2}{c}{24/50 \quad(48.0\%)}\\\midrule
Paired $t$-test $p$-value & \multicolumn{2}{c}{0.2604}\\
Conclusion & \multicolumn{2}{c}{No significant difference ($\alpha=0.05$)}\\
\bottomrule
\end{tabular}
\end{adjustbox}
\end{table}

Gemini's higher mean is driven by a few large outliers (MET:~166.7\%, BXP:~100.2\%).
The near-identical median (29.9\% vs.\ 30.8\%) and 52/48 head-to-head split confirm
that neither model dominates in expectation.

%───────────────────────────────────────────────────────────────────────────────
\section{Ensemble Methods}
%───────────────────────────────────────────────────────────────────────────────

\subsection{Methods}
\begin{enumerate}
  \item Simple Average (50/50): $\hat y=0.5F_\text{LLM}+0.5F_\text{Trad}$
  \item Median: $\hat y=\mathrm{median}(F_\text{LLM},F_\text{Trad})$
  \item Weighted 30/70: $\hat y=0.3F_\text{LLM}+0.7F_\text{Trad}$
  \item Weighted 70/30: $\hat y=0.7F_\text{LLM}+0.3F_\text{Trad}$
  \item Performance-weighted (inverse-sMAPE weights from hold-out)
  \item Hybrid / Adaptive (LLM weight = 1 for high-growth; Trad weight = 1 for stable)
\end{enumerate}

\subsection{Results}

\begin{table}[H]\centering
\caption{Ensemble performance (50 companies)}
\label{tab:ensemble}
\begin{adjustbox}{max width=\textwidth}
\begin{tabular}{lrrr}\toprule
\textbf{Method} & \textbf{Mean sMAPE} & \textbf{Median sMAPE} & \textbf{Std Dev}\\\midrule
\rowcolor{bestgreen}
\textbf{Simple Average (50/50)} & \textbf{24.14\%} & \textbf{10.73\%} & \textbf{29.04\%}\\
Median & 24.14\% & 10.73\% & 29.04\%\\
Weighted (30\% LLM, 70\% Trad) & 27.55\% & 12.34\% & 31.08\%\\
Weighted (70\% LLM, 30\% Trad) & 30.88\% & 14.67\% & 35.21\%\\
Performance-Weighted & 26.92\% & 11.89\% & 30.15\%\\
Hybrid (Adaptive) & 28.34\% & 13.02\% & 31.67\%\\\midrule
\textit{Gemini alone} & \textit{43.46\%} & \textit{29.93\%} & \textit{45.82\%}\\
\textit{Traditional alone} & \textit{36.53\%} & \textit{30.80\%} & \textit{28.47\%}\\
\bottomrule
\end{tabular}
\end{adjustbox}
\end{table}

Simple Average reduces mean sMAPE by \textbf{34\%} relative to both individual models,
at zero cost in complexity.  Weighting more toward the traditional model outperforms
equal weights in mean, but not median, indicating the traditional model is more
protective against extreme outliers.  Performance-weighting over-fits the hold-out
set and underperforms the untuned 50/50 split.

%───────────────────────────────────────────────────────────────────────────────
\section{Thinking-Mode Experiment}
\label{sec:thinking}
%───────────────────────────────────────────────────────────────────────────────

\subsection{Background}

Gemini~2.5~Flash supports a \emph{thinking} mode: before producing its final
response the model allocates a separate token budget to an internal
chain-of-thought.  Thinking tokens are invisible in the output but consume
latency and context capacity.

We test whether enabling thinking (a)~improves \emph{accuracy} and
(b)~affects \emph{robustness (consistency)}.
These are independent questions and are analysed separately below.

\subsection{Experimental Design}

\begin{table}[H]\centering
\caption{Four configurations tested (10 runs each $\times$ 3 stocks = 120 forecasts)}
\label{tab:think_configs}
\begin{adjustbox}{max width=\textwidth}
\begin{tabular}{cllrr}\toprule
\textbf{Config} & \textbf{Thinking} & \textbf{Temp} &
\textbf{Max Tokens} & \textbf{Think Budget}\\\midrule
A & Disabled & 0.0 & 8{,}192  & 0\\
B & Disabled & 0.7 & 8{,}192  & 0\\
C & Enabled  & 0.0 & 16{,}384 & 8{,}192\\
D & Enabled  & 0.7 & 16{,}384 & 8{,}192\\
\bottomrule
\end{tabular}
\end{adjustbox}
\end{table}

\textbf{Stocks:} REGN (sMAPE~4.05\%), AMZN (sMAPE~14.01\%), MET (sMAPE~166.68\%).
The thinking configs use \emph{doubled} output-token limits to compensate for
thinking-token consumption.  All 120 forecasts succeeded (100\% success rate).

\textbf{Actual 2024 values used for accuracy evaluation:}

\begin{table}[H]\centering
\caption{Actual 2024 financial data for the three test stocks}
\label{tab:actuals}
\begin{adjustbox}{max width=\textwidth}
\begin{tabular}{lrrr}\toprule
\textbf{Stock} & \textbf{Revenue} & \textbf{Net Income} & \textbf{Equity}\\\midrule
REGN & \$14.20 B & \$4.41 B  & \$29.35 B\\
AMZN & \$637.96 B & \$59.25 B & \$285.97 B\\
MET  & \$69.94 B  & \$4.43 B  & \$27.45 B\\
\bottomrule
\end{tabular}
\end{adjustbox}
\end{table}

%─────────────────────────────────────────────────────────────────────
\subsection{Analysis 1 -- Accuracy: Does Thinking Mode Improve Predictions?}
%─────────────────────────────────────────────────────────────────────

Accuracy is assessed by computing sMAPE between the \emph{mean forecast
across 10 runs} and the actual 2024 value.  Using the mean (rather than a single
run) gives each configuration the benefit of its central tendency.

\begin{table}[H]\centering
\caption{Accuracy (sMAPE of mean forecast vs.\ actual 2024 values)}
\label{tab:think_accuracy}
\begin{adjustbox}{max width=\textwidth}
\begin{tabular}{llcccr}\toprule
\textbf{Stock} & \textbf{Config} &
\textbf{Revenue sMAPE} & \textbf{Net Income sMAPE} & \textbf{Equity sMAPE} &
\textbf{Avg sMAPE}\\\midrule
\multicolumn{6}{l}{\textit{REGN (original Gemini sMAPE: 4.05\%)}}\\
& A: NoThink T=0.0 & 0.36\% & 10.14\% & 2.36\% & 4.29\%\\
& B: NoThink T=0.7 & 0.39\% &  9.76\% & 0.77\% & 3.64\%\\
& C: Think   T=0.0 & 0.34\% & 10.28\% & 0.42\% & 3.68\%\\
\rowcolor{bestgreen}
& D: Think   T=0.7 & 0.27\% &  9.46\% & 0.59\% & \textbf{3.44\%}\\[4pt]
\multicolumn{6}{l}{\textit{AMZN (original Gemini sMAPE: 14.01\%)}}\\
\rowcolor{bestgreen}
& A: NoThink T=0.0 & 0.24\% & 13.67\% & 4.76\% & \textbf{6.22\%}\\
\rowcolor{bestgreen}
& B: NoThink T=0.7 & 0.19\% &  4.68\% & 6.23\% & \textbf{3.70\%}\\
\rowcolor{badred}
& C: Think   T=0.0 & 0.30\% & 24.95\% & 4.94\% & \textbf{10.07\%}\\
\rowcolor{badred}
& D: Think   T=0.7 & 0.22\% & 25.58\% & 7.70\% & \textbf{11.17\%}\\[4pt]
\multicolumn{6}{l}{\textit{MET (original Gemini sMAPE: 166.68\%)}}\\
& A: NoThink T=0.0 & 1.81\% & 63.14\% & 4.61\% & 23.19\%\\
\rowcolor{bestgreen}
& B: NoThink T=0.7 & 1.58\% & 53.64\% & 4.70\% & \textbf{19.97\%}\\
& C: Think   T=0.0 & 1.24\% & 60.69\% & 4.93\% & 22.29\%\\
& D: Think   T=0.7 & 1.38\% & 55.00\% & 5.85\% & 20.74\%\\
\bottomrule
\end{tabular}
\end{adjustbox}
\end{table}

\textbf{Key findings -- Accuracy:}

\begin{enumerate}
  \item \textbf{Thinking mode does not improve accuracy and can worsen it.}
        For AMZN, enabling thinking at~T=0.0 more than doubles the average sMAPE
        (6.22\%~$\to$~10.07\%).  The thinking process shifts the net-income forecast
        from \$45B toward \$34--37B, moving further from the actual value of
        \$59.25B.

  \item \textbf{Temperature~0.7 slightly outperforms~0.0 in accuracy on all three
        stocks.}  Higher temperature introduces small perturbations that, on average,
        produce better-centred estimates:
        \begin{itemize}
          \item REGN: 4.29\%~(T=0.0) $\to$ 3.64\%~(T=0.7) --- $-0.65$ pp
          \item AMZN: 6.22\%~(T=0.0) $\to$ 3.70\%~(T=0.7) --- $-2.52$ pp
          \item MET:  23.19\%~(T=0.0) $\to$ 19.97\%~(T=0.7) --- $-3.22$ pp
        \end{itemize}
        This is especially pronounced for AMZN, where T=0.7 without thinking yields
        best mean net-income accuracy (\$52B vs.\ actual \$59B).

  \item \textbf{MET is uniformly poor} ($\approx$20--23\% net-income sMAPE) across all
        configurations, confirming that this stock's insurance-company accounting
        structure lies outside the model's reliable operating range.

  \item \textbf{Revenue forecasts are accurate regardless of mode} (CV~$<$~0.5\% for
        all stocks and configs), because historical revenue trend is the most
        strongly constrained input.
\end{enumerate}

%─────────────────────────────────────────────────────────────────────
\subsection{Analysis 2 -- Robustness: Does Thinking Mode Affect Consistency?}
%─────────────────────────────────────────────────────────────────────

Robustness measures how much the model's output \emph{varies across 10 independent
runs with the same prompt and parameters}.  It is independent of whether the output
is accurate.  We use Coefficient of Variation (CV, Eq.~\ref{eq:cv}).

\begin{table}[H]\centering
\caption{Robustness: Coefficient of Variation across 10 runs}
\label{tab:think_cv}
\begin{adjustbox}{max width=\textwidth}
\begin{tabular}{llcccr}\toprule
\textbf{Stock} & \textbf{Config} &
\textbf{Revenue CV} & \textbf{Net Income CV} & \textbf{Equity CV} &
\textbf{Avg CV}\\\midrule
\multicolumn{6}{l}{\textit{REGN}}\\
\rowcolor{bestgreen}
& A: NoThink T=0.0 & \textbf{0.00\%} & \textbf{0.00\%} & \textbf{0.00\%} & \textbf{0.00\%}\\
& B: NoThink T=0.7 & 0.23\%          & 3.45\%          & 3.24\%          & 2.31\%\\
& C: Think   T=0.0 & 0.26\%          & 0.48\%          & 0.87\%          & 0.54\%\\
& D: Think   T=0.7 & 0.21\%          & 3.84\%          & 0.69\%          & 1.58\%\\[4pt]
\multicolumn{6}{l}{\textit{AMZN}}\\
\rowcolor{bestgreen}
& A: NoThink T=0.0 & \textbf{0.00\%} & \textbf{0.00\%} & \textbf{0.00\%} & \textbf{0.00\%}\\
& B: NoThink T=0.7 & 0.38\%          & 8.59\%          & 2.06\%          & 3.67\%\\
& C: Think   T=0.0 & 0.14\%          & 4.46\%          &\cellcolor{badred}7.71\% & 4.10\%\\
& D: Think   T=0.7 & 0.46\%          & 1.12\%          & 6.14\%          & 2.57\%\\[4pt]
\multicolumn{6}{l}{\textit{MET}}\\
\rowcolor{bestgreen}
& A: NoThink T=0.0 & \textbf{0.00\%} & \textbf{0.00\%} & \textbf{0.00\%} & \textbf{0.00\%}\\
& B: NoThink T=0.7 & 0.34\%          & 32.41\%         & 0.61\%          & 11.12\%\\
\rowcolor{badred}
& C: Think   T=0.0 & 0.32\%          & 69.18\%         & 0.57\%          & 23.36\%\\
\rowcolor{badred}
& D: Think   T=0.7 & 0.27\%          & 75.06\%         & 3.06\%          & 26.13\%\\
\bottomrule
\end{tabular}
\end{adjustbox}
\end{table}

\textbf{Key findings -- Robustness:}

\begin{enumerate}
  \item \textbf{No Thinking + Temperature~0.0 is the only configuration that
        achieves perfectly reproducible forecasts.}
        All 10 runs produce \emph{byte-identical} JSON for every stock
        (CV~$= 0.00\%$ on every metric).  This is the only setting suitable
        for production environments requiring full auditability.

  \item \textbf{Thinking mode destroys perfect reproducibility even at T=0.0.}
        With thinking enabled, the model's internal reasoning process explores
        different chains of thought on each call (the thinking sub-network
        appears to sample at temperature $> 0$), propagating variance into the
        final numbers:
        \begin{itemize}
          \item REGN: Avg CV 0.00\%~(A) $\to$ 0.54\%~(C)
          \item AMZN: Avg CV 0.00\%~(A) $\to$ 4.10\%~(C), equity CV jumps to 7.71\%
          \item MET:  Avg CV 0.00\%~(A) $\to$ 23.36\%~(C), net-income CV 69.18\%
        \end{itemize}

  \item \textbf{Thinking mode is more damaging to robustness than raising temperature.}
        For MET, Think~T=0.0 (avg~CV~23.36\%) is worse than NoThink~T=0.7
        (avg~CV~11.12\%), which is counter-intuitive: the supposedly
        ``deterministic'' thinking mode actually introduces more variance than
        moderate stochastic sampling.

  \item \textbf{Revenue is consistently robust regardless of mode} (CV~$< 0.5\%$)
        because the historical revenue series strongly anchors the forecast;
        intermediate items (net income, equity) are more sensitive to reasoning-path
        variation.

  \item \textbf{Original sMAPE is a reliable predictor of robustness.}
        REGN (4.05\%) shows low variance everywhere; MET (166.68\%) shows high
        variance everywhere.  Companies the LLM \emph{cannot forecast accurately}
        are also the companies whose forecasts are least reproducible.
\end{enumerate}

\subsection{Latency and Token Overhead}

\begin{table}[H]\centering
\caption{Average latency and token consumption per call}
\label{tab:think_latency}
\begin{adjustbox}{max width=\textwidth}
\begin{tabular}{llccc}\toprule
\textbf{Config} & \textbf{Latency} & \textbf{Think Tokens} &
\textbf{Output Tokens} & \textbf{Throughput}\\\midrule
A: NoThink T=0.0 & 2.0 s & 0    & $\sim$309 & $\sim$30 calls/min\\
B: NoThink T=0.7 & 2.1 s & 0    & $\sim$308 & $\sim$29 calls/min\\
C: Think   T=0.0 & 20.6 s & $\sim$4{,}919 & $\sim$292 & $\sim$3 calls/min\\
D: Think   T=0.7 & 20.6 s & $\sim$4{,}826 & $\sim$301 & $\sim$3 calls/min\\
\bottomrule
\end{tabular}
\end{adjustbox}
\end{table}

Thinking adds \textbf{$\sim$5{,}000 thinking tokens per call} and increases
per-call latency by \textbf{10$\times$} (2~s~$\to$~21~s).  Throughput drops
from $\sim$30 to $\sim$3 calls/minute --- a critical constraint for batch
lending workflows.

\subsection{Consolidated Comparison}

\begin{table}[H]\centering
\caption{Summary: Accuracy vs.\ Robustness across all configurations}
\label{tab:think_summary}
\begin{adjustbox}{max width=\textwidth}
\begin{tabular}{llcccc}\toprule
\textbf{Config} & \textbf{Thinking} & \textbf{Temp} &
\textbf{Avg Accuracy} & \textbf{Avg Robustness} & \textbf{Latency}\\
& & & \textbf{(mean sMAPE)} & \textbf{(mean CV)} & \\\midrule
\rowcolor{bestgreen}
A & Disabled & 0.0 & 11.23\% & \textbf{0.00\%} & 2.0 s\\
\rowcolor{midyellow}
B & Disabled & 0.7 & \textbf{9.10\%} & 5.70\% & 2.1 s\\
\rowcolor{badred}
C & Enabled  & 0.0 & 12.01\% & 9.33\% & 20.6 s\\
\rowcolor{badred}
D & Enabled  & 0.7 & 11.12\% & 10.02\% & 20.6 s\\
\bottomrule
\end{tabular}
\end{adjustbox}
\end{table}

\textbf{Trade-off summary:}
\begin{itemize}
  \item Config~A (NoThink, T=0.0): \textbf{Best robustness (CV=0\%), near-best accuracy,
        lowest latency.} Preferred for production auditing and regulatory reporting.
  \item Config~B (NoThink, T=0.7): \textbf{Best accuracy} on average, acceptable
        robustness, identical latency.  Preferred when accuracy is the sole
        objective and multiple runs can be averaged.
  \item Configs~C and~D (With Thinking): Worse accuracy for AMZN, worse robustness
        for all stocks, 10$\times$ slower.  \textbf{Not recommended} for this task.
\end{itemize}

%───────────────────────────────────────────────────────────────────────────────
\section{Conclusions}
%───────────────────────────────────────────────────────────────────────────────

\begin{enumerate}
  \item \textbf{Gemini matches traditional forecasting in accuracy}
        ($p = 0.26$, head-to-head 52/48).  Neither model dominates.

  \item \textbf{Simple Average ensemble reduces mean sMAPE by 34\%}
        (to 24.14\%) at zero tuning cost --- the highest-impact single change.

  \item \textbf{Thinking mode does not improve accuracy and frequently worsens it.}
        For AMZN it increases avg sMAPE from 6.22\%~$\to$~10.07\% at T=0.0.
        It adds 10$\times$ latency and $\sim$5{,}000 extra tokens per call with
        no benefit.

  \item \textbf{No Thinking + Temperature 0.0 is the only perfectly reproducible
        configuration (CV~$= 0.00\%$ for all metrics on all stocks).}
        This is mandatory for regulatory auditability.

  \item \textbf{No Thinking + Temperature 0.7 provides the best mean accuracy}
        (avg sMAPE~9.10\% vs.\ 11.23\% at T=0.0) at the cost of non-deterministic
        outputs, making it suitable only when multiple runs are averaged.

  \item \textbf{Token limit must be 8{,}192 with thinking disabled.}
        4{,}096 tokens caused 40\% failure rates on large companies; adaptive
        thinking silently consumed additional budget causing partial failures
        even at 8{,}192 tokens.  Disabling thinking at 8{,}192 tokens achieves
        100\% success across all test cases.
\end{enumerate}

%───────────────────────────────────────────────────────────────────────────────
\section{Executive Recommendations for Amazon (AMZN)}
\label{sec:amazon_exec}
%───────────────────────────────────────────────────────────────────────────────

\textit{This section summarises the practical implications of our forecasting
analyses for Amazon's leadership, written for a CEO/CFO audience without
reference to modelling internals.}

\subsection*{What the Model Gets Right}

Amazon's \textbf{revenue trajectory is highly predictable}.  Across all four
experimental configurations, the AI model forecasts revenue to within
\textbf{under 0.5\% of actual 2024 values} (\$637.96 B actual vs.\
model estimates within \$3~B).  The consistent, high-growth pattern of
Amazon's top line is well-captured by trend-based methods and is robust to
any choice of model settings.

\subsection*{Where Caution is Warranted}

\textbf{Net income is materially harder to forecast.}  The best-performing
configuration (no thinking mode, temperature~0.7) estimated 2024 net income
with a 4.68\% error --- an absolute miss of approximately \$2.8~B against
the \$59.25~B actual.  Configurations using Gemini's ``thinking'' mode
performed significantly worse, producing net-income errors of 24--26\%
(roughly \$14--15~B off), demonstrating that more compute does not translate
to better financial estimates for a company of Amazon's complexity.

\subsection*{Three Actionable Recommendations}

\begin{enumerate}

  \item \textbf{Use AI-generated forecasts as a cross-check, not a sole
        source.}  The Simple Average ensemble (combining AI with a rule-based
        model) reduced mean forecast error by \textbf{34\%} compared to
        either model alone (from $\sim$14\% to $\sim$9\% for AMZN-class
        companies).  A blended approach is more reliable than relying on any
        single method for credit or strategic planning decisions.

  \item \textbf{Require deterministic outputs for audit and compliance
        workflows.}  When forecast numbers must be reproducible --- for loan
        covenants, board presentations, or regulatory filings --- the model
        must be run at temperature~0.0 with thinking disabled.  This is the
        only configuration that produces \textbf{byte-identical results on
        every run} (CV~$= 0.00\%$).  All other configurations introduce
        run-to-run variation that would be inconsistent with audit trails.

  \item \textbf{Do not activate ``deep reasoning'' (thinking) mode for
        financial forecasting.}  For AMZN specifically, enabling thinking
        \emph{more than doubled} the net-income forecast error
        (6.22\%~$\to$~10.07\% at T=0.0) while increasing API cost and
        response time by 10$\times$.  The additional reasoning did not
        provide value because Amazon's financial complexity is not reducible
        to a single chain-of-thought; the model's internal deliberation
        moved the estimate \emph{further} from the actual result.

\end{enumerate}

\subsection*{Summary Table}

\begin{table}[H]\centering
\caption{AMZN forecast quality by configuration --- key numbers for leadership}
\label{tab:amzn_exec}
\begin{adjustbox}{max width=\textwidth}
\begin{tabular}{lcccc}\toprule
\textbf{Scenario} & \textbf{Revenue error} & \textbf{Net income error} &
\textbf{Reproducible?} & \textbf{Recommended use}\\\midrule
\rowcolor{bestgreen}
AI only, no thinking, T=0.0 & 0.24\% & 13.67\% & Yes (CV=0\%) &
  Audit / compliance\\
\rowcolor{midyellow}
AI only, no thinking, T=0.7 & 0.19\% & \textbf{4.68\%}  & No (CV=3.7\%) &
  Internal planning\\
\rowcolor{bestgreen}
Ensemble (AI + traditional)  & $<$0.2\%  & $\sim$5--8\%   & Configurable &
  \textbf{Default production}\\
\rowcolor{badred}
AI with thinking, T=0.0     & 0.30\% & 24.95\% & No (CV=4.1\%) &
  Not recommended\\
\bottomrule
\end{tabular}
\end{adjustbox}
\end{table}

\subsection*{Bottom Line}

The AI forecasting system reliably tracks Amazon's revenue growth and can
serve as a cost-effective complement to analyst estimates.  For any
decision that depends on net income or earnings --- including lending
covenants, credit ratings, or capex planning --- the ensemble model
(AI~+~traditional, temperature~0.0) offers the best combination of
accuracy, reproducibility, and interpretability.  Activating advanced
``thinking'' features adds cost and latency while degrading forecast quality
for AMZN; it should remain disabled.

\end{document}